% ===========================================================================
\chapter{Reading Datalog: tokens}\label{ch:lexer}

Reading a program happens in two steps, as it does for every programming
language. The first cuts the text into \emph{tokens}, the words of the
language: a name, a variable, a number, a string, a piece of punctuation
such as \code{:-} or \code{(}. The second, in the next chapter, works out
how the tokens fit together. This first step is called \emph{lexing}.

\section{Tokens}

A token is one of a handful of kinds, so it is a union. \code{End_Of_Input}
marks the end of the program, so that the code reading tokens always has a
token to look at, even after the last real one. \code{token_text} says what
a token is, for error messages, and \code{punct_text} lets the parser ask
``is this the punctuation \code{(}?'' with a single comparison.

\begin{newfile}{lexer.nex}
-- lexer.nex: cutting Datalog text into tokens.

intern Terms

union Token
  Name_Token(text: String)
  Var_Token(text: String)
  Int_Token(value: Integer)
  Text_Token(value: String)
  Punct(text: String)
  End_Of_Input
end

function token_text(t: Token): String
do
  match t of
    Name_Token(text) then result := "'" + text + "'"
    Var_Token(text)  then result := "variable " + text
    Int_Token(value) then result := value.to_string
    Text_Token(value) then result := "\"" + value + "\""
    Punct(text)      then result := "'" + text + "'"
    End_Of_Input     then result := "the end of the program"
  end
end

-- A punctuation token's text, or "" for any other token.
function punct_text(t: Token): String
do
  match t of
    Punct(text) then result := text
    _           then result := ""
  end
end
\end{newfile}

Which kind of token starts at a character depends on what the character
is: a digit starts a number, a lower-case letter a name, a capital or an
underscore a variable.

\begin{continued}{lexer.nex}
function is_digit(c: Char): Boolean
do
  result := "0123456789".contains(c.to_string)
end

function is_lower(c: Char): Boolean
do
  result := c >= #a and c <= #z
end

function is_upper(c: Char): Boolean
do
  result := (c >= #A and c <= #Z) or c = #_
end

-- Can `c` continue a name or a variable?
function is_name_char(c: Char): Boolean
do
  result := is_lower(c) or is_upper(c) or is_digit(c)
end
\end{continued}

Nex writes a single character as \code{\#} followed by it: \code{\#a} is the
letter a. Digits are the exception, because \code{\#0} means the character
whose code is 0; so \code{is_digit} asks a string of digits instead.

\section{The lexer}

The lexer keeps the text, a position in it, and the line that position is
on, so that every error can say where it happened. Its invariant says the
position never leaves the text. \code{advance} is the only routine that
moves forward, and it is also the one place that counts lines.

\begin{continued}{lexer.nex}
class Lexer
create
  make(t: String)
  do
    text := t
    pos := 0
    line := 1
  end
feature
  text: String
  pos: Integer
  -- The line `pos` is on, counting from 1.
  line: Integer

  at_end: Boolean
  do
    result := pos >= text.length
  end

  current: Char
  require
    not_at_end: not at_end
  do
    result := text.char_at(pos)
  end

  -- Is the text at `pos` exactly `s`?
  looking_at(s: String): Boolean
  do
    result := pos + s.length <= text.length and text.substring(pos, pos + s.length) = s
  end

  advance
  require
    not_at_end: not at_end
  do
    if current = #newline then
      line := line + 1
    end
    pos := pos + 1
  ensure
    moved_on: pos = old pos + 1
  end
\end{continued}

Spaces, line breaks and comments separate tokens but are not tokens
themselves, so \code{next_token} skips them first. \code{\#space},
\code{\#newline}, \code{\#tab} and \code{\#return} are Nex's names for the
characters that cannot be written after a \code{\#} as themselves.

\begin{continued}{lexer.nex}
  -- Skip spaces, line breaks and % comments.
  skip_blanks
  do
    let done := false
    from until done do
      if at_end then
        done := true
      elseif current = #% then
        from until at_end or current = #newline do
          advance
        end
      elseif current = #space or current = #newline or current = #tab or current = #return then
        advance
      else
        done := true
      end
    end
  end

  next_token: Token
  do
    skip_blanks
    if at_end then
      result := create End_Of_Input
    elseif is_digit(current) or (current = #- and pos + 1 < text.length
                                 and is_digit(text.char_at(pos + 1))) then
      result := read_integer
    elseif is_lower(current) then
      result := create Name_Token.make(read_name)
    elseif is_upper(current) then
      result := create Var_Token.make(read_name)
    elseif current = #" then
      result := read_string
    else
      result := read_punct
    end
  ensure
    moved_on: pos > old pos or at_end
  end
\end{continued}

The postcondition of \code{next_token}, \code{moved_on}, is the reason
reading a program always ends. Every call either consumes some of the
text, or finds that there is none left and returns \code{End_Of_Input}. A
loop that calls \code{next_token} until it sees the end cannot go round
forever on the same character. Without the postcondition that argument
lives only in our heads; with it, the lexer checks its own progress on
every call.

The routines that read each kind of token come last. Each has a
precondition saying where it may start, and \code{next_token} only calls
one after looking at the character there. Punctuation is found by trying
the possible symbols in turn, longest first, so that \code{:-} is read as
one token rather than \code{:} followed by \code{-}.

\begin{continued}{lexer.nex}
  read_name: String
  require
    at_name: not at_end and is_name_char(current)
  do
    let start := pos
    from until at_end or not is_name_char(current) do
      advance
    end
    result := text.substring(start, pos)
  ensure
    not_empty: result.length > 0
  end

  read_integer: Token
  do
    let start := pos
    advance
    from until at_end or not is_digit(current) do
      advance
    end
    result := create Int_Token.make(text.substring(start, pos).to_integer)
  ensure
    moved_on: pos > old pos
  end

  read_string: Token
  require
    at_quote: not at_end and current = #"
  do
    advance
    let start := pos
    from until at_end or current = #" do
      advance
    end
    if at_end then
      raise "line " + line.to_string + ": a string is not closed"
    end
    result := create Text_Token.make(text.substring(start, pos))
    advance
  ensure
    moved_on: pos > old pos
  end

  read_punct: Token
  require
    not_at_end: not at_end
  do
    let found := ""
    across [":-", "?-", "!=", "<=", ">=", "(", ")", ",", ".", "=", "<", ">"] as p do
      if found = "" and looking_at(p) then
        found := p
      end
    end
    if found = "" then
      raise "line " + line.to_string + ": unexpected character '" + current.to_string + "'"
    end
    across found as c do
      advance
    end
    result := create Punct.make(found)
  ensure
    moved_on: pos > old pos
  end
invariant
  inside_text: 0 <= pos and pos <= text.length
  lines_count_from_one: line >= 1
end
\end{continued}

\section{Trying it out}

\begin{tryit}{try_lexer.nex}
intern Lexer

let lexer := create Lexer.make("% reach
sees(X, Z) :- follows(X, Y), sees(Y, Z).
age(\"Ann Lee\", 42).")
let done := false
from until done do
  let t := lexer.next_token
  console.print_line(lexer.line.to_string + "  " + token_text(t))
  match t of
    End_Of_Input then done := true
    _            then done := false
  end
end
\end{tryit}
\tryoutput{try_lexer}

Each token is printed with the line it ends on. The comment and the line
breaks have gone, and \code{:-} is a single token.
